% ID: SEL-D0066-OA-E02
% ESTADO: validada. S066.
\textbf{a)} En $0$, ambos límites laterales de $f$ coinciden con $f(0)=2$. Las derivadas de sus ramas son $3x^2-2x$ y $-3x^2-2x$. Como las ramas polinómicas coinciden en el punto de unión, sus pendientes laterales allí son ambas $0$. Por tanto, $f$ es continua y derivable en $0$, con $f'(0)=0$.

\textbf{b)} También para $h$ coinciden los límites laterales y el valor $2$. Las derivadas de las ramas son $-2x+1$ y $-2x-1$. Las pendientes laterales son $1$ y $-1$, así que $h$ es continua pero no derivable en $0$.

\textbf{c)} En ambos perfiles, la altura crece hasta $x=0$ y decrece después. En $h$ el cambio brusco de pendiente forma el pico del arco de la catedral. En $f$ las dos ramas se unen con tangente horizontal, formando la cima redondeada del túnel.
